\section{Diferenciación Numérica}
\subsection{Motivación}

El objetivo de esta sección es aproximar la derivada de una función $f$ en un punto $x$ a partir de los valores de la función en ciertos puntos, y llegar incluso a estimar derivadas de orden superior, es decir:

\[
    \begin{array}{c|c}
        x_i & f(x_i) \\
        \hline
        x_0 & f(x_0) \\
        x_1 & f(x_1) \\
        \vdots & \vdots \\
        x_n & f(x_n) \\
    \end{array}
    \longrightarrow
    f'(x), f''(x), \ldots f^{(m)}(x) ?
\]

Además, en este problema, se puede suponer que la partición de los puntos $x_i$ es equidistante, es decir, $x_i = x_0 + ih$ para $i = 0, 1, \ldots, n$, donde $h$ es el paso de la partición.\\

Tenemos a primera vista dos formas de abordar este problema:
\begin{itemize}
    \item Usar polinomios de interpolación para aproximar la función $f$ y luego derivar el polinomio resultante.
\end{itemize}